\section{When to Use Parallel Tempering}
\label{sec:when-to-use}

\subsection{Decision Guide}

\pt{} is the right tool when the single-chain MCMC is stuck.
The following conditions identify states where local-optimum trapping is likely:

\paragraph{Condition 1: Large $k$.}
States with $k \ge 30$ (TX, CA, FL, NY) have complex, high-dimensional plan
spaces where the \fr{} chain at $\varepsilon=0.5\%$ cannot easily traverse the
energy barriers between distinct plan regions.
The feasible set at $\varepsilon=0.5\%$ is a small fraction of the feasible
set at $\varepsilon=5\%$; the hot chain can traverse barriers in the looser
set and occasionally hand configurations to the cold chain via swaps.

\paragraph{Condition 2: Known local optima.}
If repeated runs from different seeds consistently produce the same EC value
(or a small cluster of values), the chain is likely trapped.
This is detectable from the audit log: if $\mathtt{cold\_chain\_ec\_variance}$
across multiple runs is small relative to the range of EC values seen in the
hot chain, the cold chain is trapped.

\paragraph{Condition 3: Willing to pay the replica cost.}
PT costs $N_{\mathrm{replicas}} \times$ the single-chain runtime (serial execution).
If runtime is a binding constraint (e.g., interactive use or rapid iteration),
PT is not appropriate.
For publication-quality maps or statutory submissions, the runtime cost is
usually acceptable.

\subsection{Algorithm Decision Matrix}

\begin{center}
\renewcommand{\arraystretch}{1.25}
\begin{tabular}{@{}llll@{}}
\toprule
Situation & Recommended & Reason & CLI \\
\midrule
$k \le 5$, fast iteration
  & \texttt{single}
  & Adequate mixing, minimal cost
  & \texttt{--search single} \\
$k = 6$--$20$, single best plan
  & \texttt{convergence}
  & CS sweeps seeds until convergence
  & \texttt{--search convergence} \\
$k = 6$--$20$, distributional coverage
  & \texttt{percentile}
  & Full cold chain distribution
  & \texttt{--search percentile} \\
$k \ge 30$, best plan needed
  & \textbf{PT ($N=4$)}
  & Replica exchange escapes local optima
  & \texttt{--search parallel-tempering} \\
$k \ge 30$, geographic structure strong
  & \ms{}
  & Coarsening matches geographic barriers
  & \texttt{--search multi-scale} \\
$k \ge 30$, auditable + reproducible
  & \textbf{PT ($N=4$)}
  & Exact seed derivation, full audit chain
  & \texttt{--search parallel-tempering} \\
VRA constraints present
  & VraRecom
  & VRA-aware chain preserves minority districts
  & \texttt{--search vra-recom} \\
\bottomrule
\end{tabular}
\end{center}

\subsection{PT vs.\ ForestRecom vs.\ Multi-Scale}

The three strategies for large-$k$ states are single-chain \fr{}, PT, and \ms{}.
Their trade-offs are:

\paragraph{\fr{} (baseline).}
\begin{itemize}
\item \emph{Strength}: Cheapest, simplest, exact distributional target at cold $\varepsilon$.
\item \emph{Weakness}: Stalls on TX-class states; no mechanism to escape local optima.
\item \emph{When}: Always run as baseline for comparison. Use as final algorithm
  only for $k \le 20$.
\end{itemize}

\paragraph{PT.}
\begin{itemize}
\item \emph{Strength}: No approximation. Hot chains explore broadly; cold chain
  inherits diversity. Exact \fr{} proposal distribution at every tolerance level.
  Strong audit chain (all seeds deterministic from \texttt{base\_seed}).
\item \emph{Weakness}: $N\times$ runtime cost (serial). Hot chains do not
  contribute directly to the output---they are overhead.
\item \emph{When}: $k \ge 30$, local optima suspected, runtime acceptable,
  strong auditability required.
  Default for Texas, California, Florida, New York.
\end{itemize}

\paragraph{\ms{}.}
\begin{itemize}
\item \emph{Strength}: Larger proposals---can move entire clusters of tracts
  simultaneously. Better efficiency per replica for states where the
  local-optimum barrier corresponds to a geographic boundary (e.g.,
  urban/suburban split in TX).
\item \emph{Weakness}: Coarsening introduces approximation. Requires calibrating
  the coarsening resolution $\alpha$ for each state. Complex coastlines or
  non-contiguous counties may cause coarsening failures.
\item \emph{When}: $k \ge 30$, state has strong geographic structure (river
  systems, mountain ranges) that defines the local-optimum barriers,
  and coarsening can be calibrated reliably.
\end{itemize}

\paragraph{Combining PT and \ms{}.}
In principle, each PT replica could run an \ms{} chain at its tolerance,
combining the global diversity injection of PT with the large-proposal benefit
of \ms{}.
This is not implemented in the current platform and is left as an open question
for G-series future work.

\subsection{Parameter Selection for Large States}

For large-$k$ states ($k \ge 30$), the recommended PT configuration is:
\begin{itemize}
\item $N_{\mathrm{replicas}} = 6$: Additional replicas improve ladder coverage.
  The step from $N=4$ to $N=6$ is more valuable for large states (where
  $N=4$ swap rates drop to $18\%$, below the optimal target) than for
  small states.
\item \texttt{cold\_tol} $= 0.005$: Statutory tolerance, fixed.
\item \texttt{hot\_tol} $= 0.10$: Looser than the default $0.05$ for large
  states, where the energy barrier is higher and a wider hot-chain
  sampling range is needed to bridge it.
\item \texttt{swap\_interval} $= 20$: Longer between swaps for large states,
  giving each chain more time to diffuse between swap events.
\end{itemize}
These recommendations follow from the parameter scaling table in the spec
(Section~1 of the specification document) and are supported by the
empirical swap acceptance rates in Section~\ref{sec:empirical}.
